\documentclass[11pt, a4paper]{article}
\usepackage[utf8]{inputenc}
\usepackage[margin=0.75in]{geometry}
\usepackage{amsmath, amssymb}
\usepackage{enumitem}
\usepackage{multicol}

\title{\textbf{NEET Biology: TCA Cycle (Krebs Cycle) — High-Yield Difficult Questions}}
\author{Target: NEET UG}
\date{}

\begin{document}

\maketitle

\section*{Questions}

\begin{enumerate}[label=\textbf{\arabic*.}]

\item Which of the following enzymes of the Tricarboxylic Acid (TCA) cycle is **not** located freely soluble in the mitochondrial matrix, but is instead an integral membrane protein embedded in the inner mitochondrial membrane?
\begin{enumerate}[label=(\Alph*)]
    \item Citrate synthase
    \item Isocitrate dehydrogenase
    \item Succinate dehydrogenase
    \item Malate dehydrogenase
\end{enumerate}

\item During the oxidative decarboxylation of pyruvic acid catalyzed by the Pyruvate Dehydrogenase (PDH) complex, how many coenzymes/cofactors participate in the sequential reaction mechanism?
\begin{enumerate}[label=(\Alph*)]
    \item Two ($\text{NAD}^+$ and $\text{CoA-SH}$)
    \item Three ($\text{TPP}$, $\text{NAD}^+$, and $\text{Mg}^{2+}$)
    \item Four ($\text{TPP}$, Lipoic acid, $\text{FAD}$, and $\text{NAD}^+$)
    \item Five ($\text{TPP}$, Lipoic acid, $\text{CoA-SH}$, $\text{FAD}$, and $\text{NAD}^+$)
\end{enumerate}

\item Identify the exact sequence of intermediates in the Krebs cycle starting from the condensation product of Acetyl-CoA and Oxaloacetate:
\begin{enumerate}[label=(\Alph*)]
    \item $\text{Citrate} \rightarrow \text{Isocitrate} \rightarrow \alpha\text{-Ketoglutarate} \rightarrow \text{Succinyl-CoA} \rightarrow \text{Succinate} \rightarrow \text{Fumarate} \rightarrow \text{Malate} \rightarrow \text{OAA}$
    \item $\text{Citrate} \rightarrow \alpha\text{-Ketoglutarate} \rightarrow \text{Isocitrate} \rightarrow \text{Succinyl-CoA} \rightarrow \text{Fumarate} \rightarrow \text{Succinate} \rightarrow \text{Malate} \rightarrow \text{OAA}$
    \item $\text{Isocitrate} \rightarrow \text{Citrate} \rightarrow \text{Succinyl-CoA} \rightarrow \alpha\text{-Ketoglutarate} \rightarrow \text{Succinate} \rightarrow \text{Malate} \rightarrow \text{Fumarate} \rightarrow \text{OAA}$
    \item $\text{Citrate} \rightarrow \text{Isocitrate} \rightarrow \text{Succinyl-CoA} \rightarrow \alpha\text{-Ketoglutarate} \rightarrow \text{Fumarate} \rightarrow \text{Malate} \rightarrow \text{Succinate} \rightarrow \text{OAA}$
\end{enumerate}

\item \textbf{Assertion (A):} Molecular oxygen ($\text{O}_2$) does not directly participate as a substrate in any of the eight enzymatic steps of the TCA cycle. \\
\textbf{Reason (R):} The TCA cycle ceases immediately under anaerobic conditions because $\text{NAD}^+$ and $\text{FAD}$ cannot be regenerated by the mitochondrial electron transport chain without $\text{O}_2$ as the terminal electron acceptor.
\begin{enumerate}[label=(\Alph*)]
    \item Both (A) and (R) are true and (R) is the correct explanation of (A).
    \item Both (A) and (R) are true but (R) is not the correct explanation of (A).
    \item (A) is true but (R) is false.
    \item Both (A) and (R) are false.
\end{enumerate}

\item How many molecules of $\text{H}_2\text{O}$ are consumed and how many are released, respectively, during **one complete turn** of the TCA cycle (from the entry of one Acetyl-CoA to the regeneration of Oxaloacetate, excluding the aconitase intermediate dehydration/rehydration net change)?
\begin{enumerate}[label=(\Alph*)]
    \item $2\text{ H}_2\text{O}$ consumed, $0\text{ H}_2\text{O}$ net released
    \item $1\text{ H}_2\text{O}$ consumed, $1\text{ H}_2\text{O}$ released
    \item $3\text{ H}_2\text{O}$ consumed, $1\text{ H}_2\text{O}$ released
    \item $0\text{ H}_2\text{O}$ consumed, $2\text{ H}_2\text{O}$ released
\end{enumerate}

\item Which of the following intermediates of the Krebs cycle is a **5-carbon dicarboxylic acid**?
\begin{enumerate}[label=(\Alph*)]
    \item Citric acid
    \item $\alpha$-Ketoglutaric acid
    \item Succinic acid
    \item Oxaloacetic acid
\end{enumerate}

\item Match Column-I (TCA Cycle Enzyme) with Column-II (Type of Reaction Catalyzed):
\begin{center}
\begin{tabular}{|l|l|}
\hline
\textbf{Column-I} & \textbf{Column-II} \\
\hline
P. Aconitase & 1. Substrate-level phosphorylation \\
Q. Isocitrate dehydrogenase & 2. Isomerization via dehydration-hydration \\
R. Succinyl-CoA synthetase & 3. Oxidative decarboxylation \\
S. Fumarase & 4. Hydration across a trans double bond \\
\hline
\end{tabular}
\end{center}
\begin{enumerate}[label=(\Alph*)]
    \item P-2, Q-3, R-1, S-4
    \item P-3, Q-2, R-1, S-4
    \item P-2, Q-1, R-3, S-4
    \item P-4, Q-3, R-2, S-1
\end{enumerate}

\item If Acetyl-CoA radiolabeled with $^{14}\text{C}$ at its **carbonyl carbon** ($\text{C}_1$, $\text{CH}_3-\mathbf{^{14}C}(=\text{O})-\text{SCoA}$) enters the Krebs cycle and condenses with unlabeled oxaloacetate, when will the $^{14}\text{C}$ label first be released as $^{14}\text{CO}_2$?
\begin{enumerate}[label=(\Alph*)]
    \item During the first turn of the cycle (at the isocitrate dehydrogenase step)
    \item During the first turn of the cycle (at the $\alpha$-ketoglutarate dehydrogenase step)
    \item Not at all during the first turn; $50\%$ is released during the second turn and $25\%$ during the third turn
    \item Exclusively during the oxidative decarboxylation of pyruvate
\end{enumerate}

\item Which of the following correctly lists the exact number of reducing equivalents and high-energy nucleoside triphosphates produced during **two turns** of the TCA cycle (excluding the link reaction)?
\begin{enumerate}[label=(\Alph*)]
    \item $8\text{ NADH} + \text{H}^+$, $2\text{ FADH}_2$, $2\text{ GTP}$
    \item $6\text{ NADH} + \text{H}^+$, $2\text{ FADH}_2$, $2\text{ GTP}$
    \item $6\text{ NADH} + \text{H}^+$, $4\text{ FADH}_2$, $2\text{ ATP}$
    \item $3\text{ NADH} + \text{H}^+$, $1\text{ FADH}_2$, $1\text{ GTP}$
\end{enumerate}

\item \textbf{Statement I:} Malonate is a classic competitive inhibitor of succinate dehydrogenase because it structurally resembles succinate. \\
\textbf{Statement II:} Inhibition of succinate dehydrogenase by malonate leads to the accumulation of fumarate and malate in the mitochondrial matrix.
\begin{enumerate}[label=(\Alph*)]
    \item Both Statement I and Statement II are correct.
    \item Both Statement I and Statement II are incorrect.
    \item Statement I is correct, but Statement II is incorrect.
    \item Statement I is incorrect, but Statement II is correct.
\end{enumerate}

\item Fluoroacetate (a potent rodenticide found in *Dichapetalum cymosum*) blocks the TCA cycle via ``lethal synthesis.'' How does it inhibit the cycle?
\begin{enumerate}[label=(\Alph*)]
    \item It directly inhibits pyruvate dehydrogenase by binding to lipoic acid.
    \item It is converted to fluoroacetyl-CoA, condenses with oxaloacetate to form fluorocitrate, which potently inhibits **aconitase**.
    \item It uncouples ATP synthase from the proton gradient.
    \item It competitively inhibits malate dehydrogenase.
\end{enumerate}

\item How many total $\text{ATP}$ equivalents are generated from the complete aerobic oxidation of **one molecule of pyruvic acid** inside the mitochondrion (link reaction + TCA cycle + oxidative phosphorylation)?
\begin{enumerate}[label=(\Alph*)]
    \item $12\text{ ATP}$
    \item $15\text{ ATP}$
    \item $30\text{ ATP}$
    \item $38\text{ ATP}$
\end{enumerate}

\item Which of the following steps in the TCA cycle is thermodynamically unfavorable ($\Delta G^{\circ\prime} = +29.7\text{ kJ/mol}$) in the forward direction, yet proceeds smoothly in vivo because its product is rapidly consumed by citrate synthase?
\begin{enumerate}[label=(\Alph*)]
    \item Citrate $\rightarrow$ Isocitrate
    \item $\alpha$-Ketoglutarate $\rightarrow$ Succinyl-CoA
    \item Succinate $\rightarrow$ Fumarate
    \item L-Malate $\rightarrow$ Oxaloacetate
\end{enumerate}

\item In the reaction catalyzed by **Succinyl-CoA synthetase** (succinate thiokinase), what drives the synthesis of $\text{GTP}$ (or $\text{ATP}$) from $\text{GDP} + \text{P}_i$?
\begin{enumerate}[label=(\Alph*)]
    \item Oxidation of an aldehyde group to a carboxylic acid
    \item Hydrolysis of a high-energy thioester bond in Succinyl-CoA
    \item Dehydration of a secondary alcohol
    \item Decarboxylation of a $\beta$-keto acid
\end{enumerate}

\item Which intermediate of the TCA cycle serves as the direct precursor for the biosynthesis of **porphyrins, heme, cytochromes, and chlorophyll**?
\begin{enumerate}[label=(\Alph*)]
    \item Citrate
    \item $\alpha$-Ketoglutarate
    \item Succinyl-CoA
    \item Fumarate
\end{enumerate}

\item Consider the following statements regarding the **$\alpha$-Ketoglutarate dehydrogenase complex**:
\begin{enumerate}[label=(\roman*)]
    \item It is structurally and mechanistically homologous to the pyruvate dehydrogenase complex, utilizing the same five cofactors ($\text{TPP}$, lipoate, $\text{CoA}$, $\text{FAD}$, $\text{NAD}^+$).
    \item It catalyzes the only irreversible oxidative decarboxylation step inside the TCA cycle proper.
    \item Unlike the PDH complex, it is regulated by covalent phosphorylation/dephosphorylation.
    \item Its product is a high-energy thioester compound containing $4\text{ carbons}$.
\end{enumerate}
Which of the above statements are correct?
\begin{enumerate}[label=(\Alph*)]
    \item (i), (ii), and (iv)
    \item (i) and (iii) only
    \item (ii), (iii), and (iv)
    \item (i), (ii), (iii), and (iv)
\end{enumerate}

\item Why does the TCA cycle require the continuous replenishment of its intermediates via **anaplerotic reactions**?
\begin{enumerate}[label=(\Alph*)]
    \item Because oxaloacetate is destroyed by molecular oxygen at the end of every turn.
    \item Because TCA cycle intermediates are continuously withdrawn as carbon skeletons for anabolic pathways (synthesis of amino acids, fatty acids, glucose, and nucleotides).
    \item Because citrate synthase is degraded after catalyzing $100$ turnovers.
    \item Because $\text{CO}_2$ released in the cycle depletes the carbon atoms of oxaloacetate.
\end{enumerate}

\item Which of the following is the most important **anaplerotic reaction** in mammalian liver and kidney mitochondria that directly replenishes oxaloacetate?
\begin{enumerate}[label=(\Alph*)]
    \item $\text{Pyruvate} + \text{CO}_2 + \text{ATP} + \text{H}_2\text{O} \xrightarrow{\text{Pyruvate carboxylase}} \text{Oxaloacetate} + \text{ADP} + \text{P}_i$
    \item $\text{Acetyl-CoA} + \text{CO}_2 \xrightarrow{\text{Citrate lyase}} \text{Oxaloacetate}$
    \item $\text{Succinate} + \text{CO}_2 \xrightarrow{\text{Fumarase}} \text{Malate}$
    \item $\text{Citrate} \xrightarrow{\text{Aconitase}} \text{Isocitrate}$
\end{enumerate}

\item Arsenic poisoning (specifically trivalent arsenite, $\text{AsO}_3^{3-}$) inhibits both the link reaction and the Krebs cycle by binding tightly to the vicinal dithiol ($-\text{SH}$) groups of which cofactor?
\begin{enumerate}[label=(\Alph*)]
    \item Thiamine pyrophosphate ($\text{TPP}$)
    \item Lipoic acid (lipoamide)
    \item Flavin adenine dinucleotide ($\text{FAD}$)
    \item Nicotinamide adenine dinucleotide ($\text{NAD}^+$)
\end{enumerate}

\item What is the exact stoichiometry of the net chemical equation for one complete turn of the TCA cycle?
\begin{enumerate}[label=(\Alph*)]
    \item $\text{Acetyl-CoA} + 3\text{NAD}^+ + \text{FAD} + \text{GDP} + \text{P}_i + 2\text{H}_2\text{O} \rightarrow 2\text{CO}_2 + \text{CoA-SH} + 3\text{NADH} + 3\text{H}^+ + \text{FADH}_2 + \text{GTP}$
    \item $\text{Pyruvate} + 4\text{NAD}^+ + \text{FAD} + \text{GDP} + \text{P}_i + 2\text{H}_2\text{O} \rightarrow 3\text{CO}_2 + 4\text{NADH} + 4\text{H}^+ + \text{FADH}_2 + \text{GTP}$
    \item $\text{Acetyl-CoA} + 2\text{NAD}^+ + 2\text{FAD} + \text{ATP} + \text{H}_2\text{O} \rightarrow 2\text{CO}_2 + \text{CoA-SH} + 2\text{NADH} + 2\text{FADH}_2 + \text{ADP}$
    \item $\text{Acetyl-CoA} + 3\text{NAD}^+ + \text{FAD} + \text{GDP} + \text{P}_i \rightarrow 2\text{CO}_2 + \text{CoA-SH} + 3\text{NADH} + \text{FADH}_2 + \text{GTP} + 2\text{H}_2\text{O}$
\end{enumerate}

\item Match Column-I (TCA Intermediate) with Column-II (Number of Carbons & Functional Classification):
\begin{center}
\begin{tabular}{|l|l|}
\hline
\textbf{Column-I} & \textbf{Column-II} \\
\hline
P. Citric acid & 1. $4\text{C}$ dicarboxylic acid with a trans double bond \\
Q. $\alpha$-Ketoglutaric acid & 2. $6\text{C}$ tricarboxylic acid with a tertiary hydroxyl group \\
R. Fumaric acid & 3. $4\text{C}$ keto-dicarboxylic acid \\
S. Oxaloacetic acid & 4. $5\text{C}$ keto-dicarboxylic acid \\
\hline
\end{tabular}
\end{center}
\begin{enumerate}[label=(\Alph*)]
    \item P-2, Q-4, R-1, S-3
    \item P-2, Q-3, R-1, S-4
    \item P-4, Q-2, R-3, S-1
    \item P-1, Q-4, R-2, S-3
\end{enumerate}

\item During the conversion of **Isocitrate to $\alpha$-Ketoglutarate**, an unstable enzyme-bound intermediate is transiently formed before decarboxylation occurs. Identify this intermediate:
\begin{enumerate}[label=(\Alph*)]
    \item *cis*-Aconitate
    \item Oxalosuccinate
    \item Succinyl-phosphate
    \item Enol-oxaloacetate
\end{enumerate}

\item Why does **Succinate dehydrogenase** reduce $\text{FAD}$ instead of $\text{NAD}^+$ during the oxidation of succinate to fumarate?
\begin{enumerate}[label=(\Alph*)]
    \item $\text{NAD}^+$ is absent from the mitochondrial matrix.
    \item The free energy change of oxidizing a saturated alkane ($-\text{CH}_2-\text{CH}_2-$) to an alkene ($-\text{CH}=\text{CH}-$) is insufficient to reduce $\text{NAD}^+$ ($E^{\circ\prime} = -0.32\text{ V}$), but sufficient to reduce enzyme-bound $\text{FAD}$ ($E^{\circ\prime} \approx 0.00\text{ V}$).
    \item Succinate dehydrogenase is located in the cytosol where $\text{FAD}$ predominates.
    \item Fumarate allosterically prevents $\text{NAD}^+$ binding to all dehydrogenases.
\end{enumerate}

\item Which of the following enzymes of the TCA cycle requires an **iron-sulfur ($\text{Fe}-\text{S}$) cluster** at its active site to coordinate the hydroxyl and carboxyl groups of its substrate during dehydration-hydration?
\begin{enumerate}[label=(\Alph*)]
    \item Citrate synthase
    \item Aconitase
    \item Fumarase
    \item Malate dehydrogenase
\end{enumerate}

\item \textbf{Assertion (A):} In prokaryotic cells (such as aerobic bacteria), the enzymes of the TCA cycle are located in the cytosol, except succinate dehydrogenase which is bound to the plasma membrane. \\
\textbf{Reason (R):} Prokaryotes lack membrane-bound organelles such as mitochondria, so their plasma membrane performs the bioenergetic functions of the inner mitochondrial membrane.
\begin{enumerate}[label=(\Alph*)]
    \item Both (A) and (R) are true and (R) is the correct explanation of (A).
    \item Both (A) and (R) are true but (R) is not the correct explanation of (A).
    \item (A) is true but (R) is false.
    \item Both (A) and (R) are false.
\end{enumerate}

\item In plants and microorganisms (but **not** in vertebrates), the **Glyoxylate Cycle** bypasses the two decarboxylation steps of the TCA cycle, allowing the net synthesis of carbohydrates from stored fatty acids. Which two unique enzymes enable this bypass?
\begin{enumerate}[label=(\Alph*)]
    \item Isocitrate lyase and Malate synthase
    \item Citrate lyase and Pyruvate carboxylase
    \item Aconitase and Fumarase
    \item Succinate thiokinase and PEP carboxykinase
\end{enumerate}

\item How many **substrate-level phosphorylation** events occur during the complete oxidation of **one molecule of glucose** through glycolysis and the TCA cycle combined?
\begin{enumerate}[label=(\Alph*)]
    \item 2
    \item 4
    \item 6
    \item 38
\end{enumerate}

\item Which three enzymes serve as the primary regulatory control points (allosteric pacemakers) of the TCA cycle proper?
\begin{enumerate}[label=(\Alph*)]
    \item Aconitase, Fumarase, Malate dehydrogenase
    \item Citrate synthase, Isocitrate dehydrogenase, $\alpha$-Ketoglutarate dehydrogenase
    \item Succinyl-CoA synthetase, Succinate dehydrogenase, Fumarase
    \item Hexokinase, Phosphofructokinase-1, Pyruvate kinase
\end{enumerate}

\item In eukaryotic cells, **citrate** can exit the mitochondrial matrix into the cytosol via the tricarboxylate carrier when $\text{ATP}$ and $\text{NADH}$ are abundant. In the cytosol, citrate serves two crucial roles:
\begin{enumerate}[label=(\Alph*)]
    \item It activates phosphofructokinase-1 ($\text{PFK-1}$) and inhibits acetyl-CoA carboxylase.
    \item It inhibits phosphofructokinase-1 ($\text{PFK-1}$) to slow glycolysis and provides Acetyl-CoA (via ATP-citrate lyase) for **fatty acid synthesis**.
    \item It is converted directly into glucose-6-phosphate by cytosolic aconitase.
    \item It stimulates lactate dehydrogenase to initiate fermentation.
\end{enumerate}

\item What is the total number of $\text{CO}_2$ molecules released when **one molecule of sucrose** undergoes complete aerobic respiration (glycolysis + link reaction + TCA cycle)?
\begin{enumerate}[label=(\Alph*)]
    \item 4
    \item 6
    \item 8
    \item 12
\end{enumerate}

\item \textbf{Statement I:} During the conversion of succinyl-CoA to succinate, a molecule of $\text{GTP}$ is synthesized in animal tissues, whereas $\text{ATP}$ is directly synthesized in plants and bacteria. \\
\textbf{Statement II:} In animal mitochondria, $\text{GTP}$ transfers its terminal phosphate group to $\text{ADP}$ via the enzyme nucleoside diphosphate kinase ($\text{GTP} + \text{ADP} \leftrightarrow \text{GDP} + \text{ATP}$).
\begin{enumerate}[label=(\Alph*)]
    \item Both Statement I and Statement II are correct.
    \item Both Statement I and Statement II are incorrect.
    \item Statement I is correct, but Statement II is incorrect.
    \item Statement I is incorrect, but Statement II is correct.
\end{enumerate}

\item How many redox reactions occur inside the **TCA cycle proper** (from citrate to oxaloacetate) per turn of the cycle?
\begin{enumerate}[label=(\Alph*)]
    \item 2
    \item 3
    \item 4
    \item 5
\end{enumerate}

\item Which of the following intermediates of the Krebs cycle is optically active (chiral) and exists specifically as the **L-stereoisomer** in the cycle?
\begin{enumerate}[label=(\Alph*)]
    \item Citrate
    \item Succinate
    \item Fumarate
    \item Malate
\end{enumerate}

\item If a student incubates isolated, intact mitochondria with **pyruvate, $\text{ADP}$, and $\text{P}_i$** in an aerobic buffer, respiration proceeds rapidly. However, if **malonate** is added, oxygen consumption stops. Adding a small catalytic amount of which compound will restore oxygen consumption for one single oxidation step even in the presence of malonate?
\begin{enumerate}[label=(\Alph*)]
    \item Succinate
    \item Fumarate or Malate
    \item Glucose
    \item Lactate
\end{enumerate}

\item Which amino acids are directly synthesized via transamination of **$\alpha$-ketoglutarate** and **oxaloacetate**, respectively?
\begin{enumerate}[label=(\Alph*)]
    \item Alanine and Glycine
    \item Glutamate and Aspartate
    \item Serine and Cysteine
    \item Valine and Leucine
\end{enumerate}

\item During the link reaction (oxidative decarboxylation of pyruvate), which carbon of pyruvic acid is released as $\text{CO}_2$?
\begin{enumerate}[label=(\Alph*)]
    \item The carboxyl carbon ($\text{C}_1$, $-\text{COOH}$)
    \item The carbonyl carbon ($\text{C}_2$, $-\text{C}=\text{O}$)
    \item The methyl carbon ($\text{C}_3$, $-\text{CH}_3$)
    \item Any of the three carbons randomly
\end{enumerate}

\item Why is the TCA cycle described as an **amphibolic pathway** rather than purely catabolic?
\begin{enumerate}[label=(\Alph*)]
    \item Because it operates simultaneously in both the cytosol and the mitochondrial matrix.
    \item Because it functions in both aerobic and anaerobic organisms identically.
    \item Because it involves both catabolic breakdown of acetyl groups to $\text{CO}_2$ for energy and anabolic withdrawal of intermediates for biosynthesis of biomolecules.
    \item Because it produces both $\text{ATP}$ and $\text{GTP}$ in equal proportions.
\end{enumerate}

\item In the reaction catalyzed by **Citrate synthase**, Acetyl-CoA ($2\text{C}$) condenses with Oxaloacetate ($4\text{C}$) to form Citrate ($6\text{C}$). What is the transient enzyme-bound intermediate formed before hydrolysis releases $\text{CoA-SH}$?
\begin{enumerate}[label=(\Alph*)]
    \item *cis*-Aconitate
    \item Citroyl-CoA
    \item Oxalosuccinate
    \item Malonyl-CoA
\end{enumerate}

\item Which of the following statements about the **Pyruvate Dehydrogenase (PDH) complex** is **false**?
\begin{enumerate}[label=(\Alph*)]
    \item It is located in the mitochondrial matrix in eukaryotes.
    \item It is activated by high ratios of $\text{NAD}^+/\text{NADH}$, $\text{CoA}/\text{Acetyl-CoA}$, and $\text{ADP}/\text{ATP}$.
    \item In mammals, phosphorylation of the $\text{E}_1$ subunit by PDH kinase **activates** the enzyme complex.
    \item Calcium ions ($\text{Ca}^{2+}$) released during muscle contraction stimulate PDH phosphatase, thereby activating the PDH complex.
\end{enumerate}

\item Calculate the total number of $\text{ATP}$ molecules synthesized via **oxidative phosphorylation alone** (excluding substrate-level phosphorylation) from the reducing equivalents generated during **one turn** of the TCA cycle (excluding the link reaction):
\begin{enumerate}[label=(\Alph*)]
    \item $9\text{ ATP}$
    \item $11\text{ ATP}$
    \item $12\text{ ATP}$
    \item $14\text{ ATP}$
\end{enumerate}

\item Which of the following enzymes in the TCA cycle is **strictly dependent on $\text{Mn}^{2+}$ or $\text{Mg}^{2+}$** to stabilize the enolate/oxalosuccinate intermediate during decarboxylation?
\begin{enumerate}[label=(\Alph*)]
    \item Fumarase
    \item Isocitrate dehydrogenase
    \item Malate dehydrogenase
    \item Citrate synthase
\end{enumerate}

\item What is the fate of the two carbon atoms of the **acetyl group** of Acetyl-CoA during the **first turn** of the Krebs cycle?
\begin{enumerate}[label=(\Alph*)]
    \item Both carbons are immediately lost as $2\text{ CO}_2$ molecules during the first turn.
    \item Neither carbon is lost as $\text{CO}_2$ during the first turn; the two $\text{CO}_2$ molecules released come from the oxaloacetate molecule that condensed with Acetyl-CoA.
    \item One carbon is lost at isocitrate dehydrogenase and the other at $\alpha$-ketoglutarate dehydrogenase.
    \item Both carbons remain permanently trapped in succinate and never leave the cycle.
\end{enumerate}

\item In **Fumarase** (fumarate hydratase), the hydration of fumarate is strictly stereospecific. Why does fumarase act on **fumarate** (*trans*-butenedioate) and **not** on **maleate** (*cis*-butenedioate)?
\begin{enumerate}[label=(\Alph*)]
    \item Maleate is a $6\text{-carbon}$ tricarboxylic acid.
    \item Fumarase possesses a stereospecific active site that only binds the *trans* double bond configuration, adding $-\text{OH}$ and $-\text{H}$ exclusively in *trans* orientation to form L-malate.
    \item Maleate is instantly reduced to ethanol in the matrix.
    \item Fumarase requires $\text{NADH}$ to bind *cis* isomers.
\end{enumerate}

\item \textbf{Statement I:} Krebs cycle begins with the condensation of acetyl group with oxaloacetic acid and water to yield citric acid. \\
\textbf{Statement II:} During the conversion of succinyl-CoA to succinic acid, a molecule of $\text{GTP}$ is synthesized, which is an example of oxidative phosphorylation.
\begin{enumerate}[label=(\Alph*)]
    \item Both Statement I and Statement II are correct.
    \item Both Statement I and Statement II are incorrect.
    \item Statement I is correct, but Statement II is incorrect.
    \item Statement I is incorrect, but Statement II is correct.
\end{enumerate}

\item If $3\text{ molecules}$ of **acetyl-CoA** enter the Krebs cycle, how many total molecules of $\text{CO}_2$, $\text{NADH} + \text{H}^+$, $\text{FADH}_2$, and $\text{GTP}$ are produced, respectively?
\begin{enumerate}[label=(\Alph*)]
    \item $6\text{ CO}_2$, $9\text{ NADH} + \text{H}^+$, $3\text{ FADH}_2$, $3\text{ GTP}$
    \item $3\text{ CO}_2$, $9\text{ NADH} + \text{H}^+$, $3\text{ FADH}_2$, $3\text{ GTP}$
    \item $6\text{ CO}_2$, $12\text{ NADH} + \text{H}^+$, $3\text{ FADH}_2$, $3\text{ GTP}$
    \item $9\text{ CO}_2$, $9\text{ NADH} + \text{H}^+$, $6\text{ FADH}_2$, $3\text{ GTP}$
\end{enumerate}

\end{enumerate}

\newpage

\section*{Answer Key}

\begin{multicols}{5}
\begin{enumerate}[label=\textbf{\arabic*.}]
    \item (C)
    \item (D)
    \item (A)
    \item (A)
    \item (A)
    \item (B)
    \item (A)
    \item (C)
    \item (B)
    \item (C)
    \item (B)
    \item (B)
    \item (D)
    \item (B)
    \item (C)
    \item (A)
    \item (B)
    \item (A)
    \item (B)
    \item (A)
    \item (A)
    \item (B)
    \item (B)
    \item (B)
    \item (A)
    \item (A)
    \item (C)
    \item (B)
    \item (B)
    \item (D)
    \item (A)
    \item (C)
    \item (D)
    \item (B)
    \item (B)
    \item (A)
    \item (C)
    \item (B)
    \item (C)
    \item (B)
    \item (B)
    \item (B)
    \item (B)
    \item (C)
    \item (A)
\end{enumerate}
\end{multicols}

\section*{Detailed Solutions}

\begin{enumerate}[label=\textbf{\arabic*.}]

\item \textbf{(C)} **Succinate dehydrogenase** is the only enzyme of the TCA cycle embedded in the inner mitochondrial membrane (where it also functions as Complex II of the electron transport chain). All other seven enzymes are soluble in the mitochondrial matrix.

\item \textbf{(D)} As explicitly stated in NCERT Class 11 Biology: \textit{``The reaction is catalysed by pyruvic dehydrogenase, which requires the participation of several coenzymes, including $\text{NAD}^+$ and Coenzyme A''}, along with $\text{TPP}$, lipoic acid, and $\text{FAD}$ — making **five cofactors/coenzymes** in total ($\text{TPP}$, Lipoic acid, $\text{CoA-SH}$, $\text{FAD}$, $\text{NAD}^+$).

\item \textbf{(A)} The canonical sequence is: Citrate ($6\text{C}$) $\rightarrow$ *cis*-Aconitate $\rightarrow$ Isocitrate ($6\text{C}$) $\rightarrow$ $\alpha$-Ketoglutarate ($5\text{C}$) $\rightarrow$ Succinyl-CoA ($4\text{C}$) $\rightarrow$ Succinate ($4\text{C}$) $\rightarrow$ Fumarate ($4\text{C}$) $\rightarrow$ Malate ($4\text{C}$) $\rightarrow$ Oxaloacetate ($4\text{C}$).

\item \textbf{(A)} Although $\text{O}_2$ is not a direct reactant in the TCA cycle, the cycle strictly depends on $\text{O}_2$ because $\text{NAD}^+$ and $\text{FAD}$ can only be regenerated by passing electrons down the ETC to $\text{O}_2$. Without $\text{O}_2$, $\text{NAD}^+$ and $\text{FAD}$ are depleted, halting the cycle.

\item \textbf{(A)} Two molecules of $\text{H}_2\text{O}$ are net consumed per turn: $1\text{ H}_2\text{O}$ at **Citrate synthase** (to hydrolyze citroyl-CoA) and $1\text{ H}_2\text{O}$ at **Fumarase** (to hydrate fumarate to malate). Aconitase removes $1\text{ H}_2\text{O}$ and adds $1\text{ H}_2\text{O}$ back, giving $0$ net release.

\item \textbf{(B)} Citric acid is $6\text{C}$ (tricarboxylic); **$\alpha$-Ketoglutaric acid** is $5\text{C}$ (dicarboxylic); Succinic acid and Oxaloacetic acid are $4\text{C}$ (dicarboxylic).

\item \textbf{(A)} Aconitase isomerizes citrate to isocitrate via dehydration-hydration (P-2). Isocitrate dehydrogenase catalyzes oxidative decarboxylation (Q-3). Succinyl-CoA synthetase performs substrate-level phosphorylation to yield $\text{GTP}/\text{ATP}$ (R-1). Fumarase hydrates fumarate across its trans double bond (S-4).

\item \textbf{(C)} During the first turn, the two $\text{CO}_2$ molecules lost come from the oxaloacetate carbons, not the incoming acetyl carbons. Because succinate is a symmetric molecule, the carbonyl $^{14}\text{C}$ label becomes randomized equally between $\text{C}_1$ and $\text{C}_4$ of oxaloacetate at the end of Turn 1. Thus, **$50\%$** is lost as $^{14}\text{CO}_2$ in Turn 2, $25\%$ in Turn 3, etc.

\item \textbf{(B)} One turn of the TCA cycle produces $3\text{ NADH} + \text{H}^+$, $1\text{ FADH}_2$, and $1\text{ GTP}$ (or $\text{ATP}$). Therefore, two turns (corresponding to $1\text{ glucose}$) produce **$6\text{ NADH} + \text{H}^+$**, **$2\text{ FADH}_2$**, and **$2\text{ GTP}$**.

\item \textbf{(C)} Statement I is true: malonate ($\text{HOOC}-\text{CH}_2-\text{COOH}$) mimics succinate ($\text{HOOC}-\text{CH}_2-\text{CH}_2-\text{COOH}$) and competitively inhibits succinate dehydrogenase. Statement II is false: blocking succinate dehydrogenase causes its substrate (**succinate**) to accumulate, while downstream products (fumarate, malate) become depleted.

\item \textbf{(B)} Fluoroacetate is activated to fluoroacetyl-CoA and condensed with oxaloacetate by citrate synthase to form **fluorocitrate**, a potent suicide inhibitor of **aconitase**.

\item \textbf{(B)} Complete oxidation of $1\text{ pyruvate}$ yields:
\begin{itemize}
    \item Link reaction: $1\text{ NADH} = 3\text{ ATP}$
    \item TCA cycle: $3\text{ NADH} = 9\text{ ATP}$, $1\text{ FADH}_2 = 2\text{ ATP}$, $1\text{ GTP} = 1\text{ ATP}$
\end{itemize}
Total = $3 + 9 + 2 + 1 = \mathbf{15\text{ ATP}}$ (Note: $2\text{ pyruvates}$ from $1\text{ glucose}$ give $30\text{ ATP}$, plus $8\text{ ATP}$ from glycolysis = $38\text{ ATP}$).

\item \textbf{(D)} The oxidation of L-malate to oxaloacetate by **malate dehydrogenase** is strongly endergonic ($\Delta G^{\circ\prime} = +29.7\text{ kJ/mol}$), keeping intracellular oxaloacetate concentrations very low. It is pulled forward by the highly exergonic citrate synthase reaction ($\Delta G^{\circ\prime} = -31.5\text{ kJ/mol}$).

\item \textbf{(B)} Succinyl-CoA contains a high-energy **thioester bond** ($\Delta G^{\circ\prime} \approx -36\text{ kJ/mol}$). Cleavage of this thioester bond by succinyl-CoA synthetase releases enough free energy to drive the phosphorylation of $\text{GDP}$ (or $\text{ADP}$) to $\text{GTP}$ (or $\text{ATP}$).

\item \textbf{(C)} **Succinyl-CoA** condenses with glycine (via $\delta$-aminolevulinic acid synthase) to initiate the synthesis of porphyrins, heme, cytochromes, and chlorophyll.

\item \textbf{(A)} Statements (i), (ii), and (iv) are correct. Statement (iii) is incorrect because unlike the PDH complex, the $\alpha$-ketoglutarate dehydrogenase complex is **not** regulated by covalent phosphorylation/dephosphorylation; it is regulated allosterically by $\text{Ca}^{2+}$, $\text{ATP}$, $\text{NADH}$, and succinyl-CoA.

\item \textbf{(B)} Because intermediates like $\alpha$-ketoglutarate, succinyl-CoA, and oxaloacetate are constantly drained off to synthesize amino acids, nucleotides, porphyrins, and glucose, **anaplerotic (filling-up) reactions** are required to replenish oxaloacetate and keep the cycle running.

\item \textbf{(A)} **Pyruvate carboxylase** (a biotin-dependent mitochondrial enzyme activated by Acetyl-CoA) carboxylates pyruvate to directly replenish oxaloacetate.

\item \textbf{(B)} Trivalent arsenite ($\text{AsO}_3^{3-}$) and mercury ($\text{Hg}^{2+}$) form stable cyclic adducts with the two adjacent thiol ($-\text{SH}$) groups of reduced **lipoic acid (dihydrolipoamide)**, completely paralyzing both PDH and $\alpha$-ketoglutarate dehydrogenase.

\item \textbf{(A)} The exact balanced equation for one turn of the TCA cycle (from Acetyl-CoA) is:
$$\text{Acetyl-CoA} + 3\text{NAD}^+ + \text{FAD} + \text{GDP} + \text{P}_i + 2\text{H}_2\text{O} \rightarrow 2\text{CO}_2 + \text{CoA-SH} + 3\text{NADH} + 3\text{H}^+ + \text{FADH}_2 + \text{GTP}$$

\item \textbf{(A)} Citric acid is a $6\text{C}$ tricarboxylic acid (P-2). $\alpha$-Ketoglutaric acid is a $5\text{C}$ keto-dicarboxylic acid (Q-4). Fumaric acid is a $4\text{C}$ dicarboxylic acid with a *trans* double bond (R-1). Oxaloacetic acid is a $4\text{C}$ keto-dicarboxylic acid (S-3).

\item \textbf{(B)} Isocitrate dehydrogenase first oxidizes the secondary alcohol of isocitrate to a ketone, forming the unstable enzyme-bound $\beta$-keto acid intermediate **oxalosuccinate**, which immediately undergoes $\text{Mn}^{2+}$-assisted decarboxylation to $\alpha$-ketoglutarate.

\item \textbf{(B)} Oxidizing a saturated $-\text{CH}_2-\text{CH}_2-$ bond to an unsaturated $-\text{CH}=\text{CH}-$ bond has a standard reduction potential of $+0.03\text{ V}$, which cannot reduce $\text{NAD}^+$ ($-0.32\text{ V}$). Instead, electrons are transferred to covalently bound **$\text{FAD}$** ($E^{\circ\prime} \approx 0.00\text{ V}$).

\item \textbf{(B)} **Aconitase** is an iron-sulfur protein containing a $[4\text{Fe}-4\text{S}]$ cluster that directly coordinates the substrate hydroxyl and carboxyl groups (rather than functioning in electron transport).

\item \textbf{(A)} Both statements are true: prokaryotes lack mitochondria, so soluble TCA enzymes reside in the cytosol while membrane-bound succinate dehydrogenase (Complex II) and ETC complexes reside in the plasma membrane.

\item \textbf{(A)} Plants (especially germinating oilseeds) and bacteria possess **isocitrate lyase** (cleaving isocitrate into succinate and glyoxylate) and **malate synthase** (condensing glyoxylate with Acetyl-CoA to form malate), enabling net conversion of fats into carbohydrates.

\item \textbf{(C)} Per glucose molecule: glycolysis generates **$4\text{ ATP}$** via substrate-level phosphorylation ($2$ at phosphoglycerate kinase + $2$ at pyruvate kinase), and two turns of the TCA cycle generate **$2\text{ GTP}/\text{ATP}$** at succinyl-CoA synthetase. Total = $4 + 2 = \mathbf{6\text{ substrate-level phosphorylations}}$.

\item \textbf{(B)} **Citrate synthase**, **Isocitrate dehydrogenase**, and **$\alpha$-Ketoglutarate dehydrogenase** catalyze the three physiologically irreversible, strongly exergonic steps of the TCA cycle and serve as its primary allosteric regulatory points.

\item \textbf{(B)} High cytosolic **citrate** signals energy abundance: it allosterically inhibits $\text{PFK-1}$ (slowing glycolysis) and is cleaved by ATP-citrate lyase into oxaloacetate and **Acetyl-CoA** to fuel **de novo fatty acid synthesis**.

\item \textbf{(D)} One molecule of sucrose ($\text{C}_{12}\text{H}_{22}\text{O}_{11}$) is hydrolyzed into $1\text{ glucose} + 1\text{ fructose}$ ($2\text{ hexoses} = 12\text{ carbons}$). Complete aerobic oxidation releases all $12\text{ carbons}$ as **$12\text{ molecules}$ of $\text{CO}_2$** ($4$ in link reaction + $8$ in TCA cycle).

\item \textbf{(A)} Both statements are correct. Animals possess a GDP-specific succinyl-CoA synthetase (followed by nucleoside diphosphate kinase to convert $\text{GTP} \rightarrow \text{ATP}$), whereas plants and bacteria possess an ADP-specific enzyme that directly forms $\text{ATP}$.

\item \textbf{(C)} **Four** redox reactions occur per turn of the TCA cycle:
\begin{enumerate}
    \item Isocitrate $\rightarrow$ $\alpha$-Ketoglutarate ($\text{NAD}^+ \rightarrow \text{NADH}$)
    \item $\alpha$-Ketoglutarate $\rightarrow$ Succinyl-CoA ($\text{NAD}^+ \rightarrow \text{NADH}$)
    \item Succinate $\rightarrow$ Fumarate ($\text{FAD} \rightarrow \text{FADH}_2$)
    \item Malate $\rightarrow$ Oxaloacetate ($\text{NAD}^+ \rightarrow \text{NADH}$)
\end{enumerate}

\item \textbf{(D)} **L-Malate** (and Isocitrate) are chiral molecules. Citrate, succinate, fumarate, $\alpha$-ketoglutarate, succinyl-CoA, and oxaloacetate are achiral (symmetrical/planar).

\item \textbf{(B)} Malonate blocks succinate dehydrogenase, preventing oxaloacetate regeneration. Adding **fumarate or malate** bypasses the block, allowing malate dehydrogenase to oxidize malate to oxaloacetate and generate $1\text{ NADH}$ (which reduces $\text{O}_2$ via the ETC).

\item \textbf{(B)} Transamination of **$\alpha$-ketoglutarate** yields **glutamate**, and transamination of **oxaloacetate** yields **aspartate**.

\item \textbf{(A)} Pyruvate dehydrogenase decarboxylates the **carboxyl group ($\text{C}_1$, $-\text{COOH}$)** of pyruvate as $\text{CO}_2$, leaving the $\text{C}_2$ carbonyl and $\text{C}_3$ methyl carbons to form the acetyl group of Acetyl-CoA.

\item \textbf{(C)} As highlighted in NCERT, the respiratory pathway (especially the TCA cycle) is an **amphibolic pathway** because it participates in both **catabolism** (breaking down substrates for energy) and **anabolism** (supplying carbon skeletons for synthesizing fatty acids, amino acids, etc.).

\item \textbf{(B)} Citrate synthase condenses Acetyl-CoA and oxaloacetate to form a transient enzyme-bound thioester intermediate, **citroyl-CoA**, whose subsequent hydrolysis drives the reaction strongly forward.

\item \textbf{(C)} Statement (C) is false because in mammals, phosphorylation of the $\text{E}_1$ subunit of the PDH complex by PDH kinase **inactivates** (inhibits) the enzyme, whereas dephosphorylation by PDH phosphatase **activates** it.

\item \textbf{(B)} One turn of the TCA cycle produces $3\text{ NADH} + \text{H}^+$ and $1\text{ FADH}_2$. Via oxidative phosphorylation: $(3 \times 3\text{ ATP}) + (1 \times 2\text{ ATP}) = 9 + 2 = \mathbf{11\text{ ATP}}$ (excluding the $1\text{ GTP}$ formed by substrate-level phosphorylation).

\item \textbf{(B)} **Isocitrate dehydrogenase** requires divalent metal ions ($\text{Mn}^{2+}$ or $\text{Mg}^{2+}$) to coordinate the carbonyl oxygen of the transient oxalosuccinate intermediate and facilitate decarboxylation.

\item \textbf{(B)} During the first turn of the TCA cycle, **neither** of the two carbons entering via Acetyl-CoA is lost as $\text{CO}_2$; both $\text{CO}_2$ molecules released come from the oxaloacetate molecule that accepted the acetyl group.

\item \textbf{(B)} Fumarase exhibits strict geometric and stereospecificity: it binds only the *trans* isomer (**fumarate**) and never the *cis* isomer (maleate), hydrating it exclusively to **L-malate**.

\item \textbf{(C)} Statement I is correct (verbatim NCERT). Statement II is incorrect because the synthesis of $\text{GTP}$ during the conversion of succinyl-CoA to succinic acid is an example of **substrate-level phosphorylation**, not oxidative phosphorylation.

\item \textbf{(A)} Per Acetyl-CoA: $2\text{ CO}_2$, $3\text{ NADH} + \text{H}^+$, $1\text{ FADH}_2$, and $1\text{ GTP}$ are produced. Multiplying by $3\text{ molecules}$ of Acetyl-CoA gives **$6\text{ CO}_2$, $9\text{ NADH} + \text{H}^+$, $3\text{ FADH}_2$, and $3\text{ GTP}$**.

\end{enumerate}

\end{document}
</USER_REQUEST>
<ADDITIONAL_METADATA>
The current local time is: 2026-11-16T12:01:21+05:30.

The user's current state is as follows:
Active Document: /Users/laxmikumari/Desktop/web/testseries/scripts/add_botany_zoology_to_part_11.mjs (LANGUAGE_UNSPECIFIED)
Cursor is on line: 1
Other open documents:
- /Users/laxmikumari/Desktop/web/testseries/.env.local (LANGUAGE_UNSPECIFIED)
- /Users/laxmikumari/Downloads/neet_five_kingdom_classification_test.tex (LANGUAGE_UNSPECIFIED)
- /Users/laxmikumari/Downloads/GOC_Extracted_Questions.tex (LANGUAGE_UNSPECIFIED)
- /Users/laxmikumari/Downloads/neet_biology_question_bank_phases_of_growth_set_2.tex (LANGUAGE_UNSPECIFIED)
- /Users/laxmikumari/Desktop/web/testseries/src/app/api/admin/notifications/route.js (LANGUAGE_UNSPECIFIED)
</ADDITIONAL_METADATA>